\chapter{Password Hashing with bcrypt}

A database full of user records is one breach away from being public. If
passwords are stored as plain text, that breach instantly compromises every
account -- and because people reuse passwords, it compromises their accounts
on other sites too. This chapter covers the one non-negotiable rule of
authentication: passwords are never stored as-is.

\begin{warnbox}[WARNING]
Never store passwords in plain text, and never store them encrypted with a
reversible cipher either. Encryption implies a key that can decrypt the
value back to the original password -- if that key leaks, every password
leaks with it. Hashing is one-way: there is no key that turns a hash back
into the original password. This is why bcrypt, not AES or a custom cipher,
is the right tool for passwords.
\end{warnbox}

\section{Hashing and Salt}

A hash function takes an input of any length and produces a fixed-length
output. The same input always produces the same output, but there is no
practical way to reverse it. If two users pick the password
\texttt{password123}, a plain hash function (like SHA-256) would produce the
\emph{same} hash for both -- which lets an attacker use precomputed lookup
tables (``rainbow tables'') to crack thousands of hashes at once.

A \textbf{salt} fixes this. It is random data mixed into the password before
hashing, so the same password produces a different hash every time. bcrypt
generates and stores this salt automatically as part of the hash string
itself, so you never manage salts manually.

\begin{diagrambox}[bcrypt Hashing Flow]
\begin{verbatim}
 Plain password: "MySecret123"
        |
        v
 bcrypt generates a random salt
        |
        v
 Salt + password run through the bcrypt algorithm
 (repeated for 2^rounds iterations)
        |
        v
 Resulting hash stored in DB:
 $2b$10$N9qo8uLOickgx2ZMRZoMy...
 (algorithm + rounds + salt + hash, all in one string)
\end{verbatim}
\end{diagrambox}

\section{Hashing on Registration}

Install bcrypt and hash the password before saving the user document.

\begin{lstlisting}[language=JavaScript]
// controllers/authController.js
import bcrypt from "bcrypt";
import User from "../models/User.js";

export const registerUser = async (req, res, next) => {
  try {
    const { name, email, password } = req.body;

    const existing = await User.findOne({ email });
    if (existing) {
      return res.status(409).json({ success: false, message: "Email already in use" });
    }

    const salt = 10; // salt rounds
    const hashedPassword = await bcrypt.hash(password, salt);

    const user = await User.create({
      name,
      email,
      password: hashedPassword,
      role: "user",
    });

    res.status(201).json({
      success: true,
      data: { id: user._id, name: user.name, email: user.email, role: user.role },
    });
  } catch (err) {
    next(err);
  }
};
\end{lstlisting}

\subsection{The Salt Rounds Tradeoff}

The second argument to \texttt{bcrypt.hash} is the \textbf{cost factor}
(commonly called salt rounds). It controls how many times the hashing
algorithm loops internally -- work scales as $2^{\text{rounds}}$.

\begin{center}
\begin{tabularx}{\textwidth}{lXX}
\toprule
\textbf{Rounds} & \textbf{Effect} & \textbf{Tradeoff} \\
\midrule
8--10 & Fast (a few ms to ~100ms per hash) & Weaker against brute force, fine for most apps \\
12 & Slower (~250ms per hash) & Better resistance, noticeable login latency \\
14+ & Very slow (seconds) & Strong resistance, but hurts login UX and server throughput \\
\bottomrule
\end{tabularx}
\end{center}

\begin{notebox}[NOTE]
10 is the widely used default and a reasonable choice for a typical Node.js
API -- it costs bcrypt roughly 50--100ms per hash on modern hardware, which
is imperceptible to a user logging in but expensive enough to make brute
forcing millions of guesses impractical.
\end{notebox}

\section{Verifying on Login}

You never ``unhash'' a password to compare it. Instead, bcrypt re-hashes the
submitted password using the salt embedded in the stored hash, and compares
the results.

\begin{lstlisting}[language=JavaScript]
// controllers/authController.js
export const loginUser = async (req, res, next) => {
  try {
    const { email, password } = req.body;

    const user = await User.findOne({ email });
    if (!user) {
      return res.status(401).json({ success: false, message: "Invalid credentials" });
    }

    const match = await bcrypt.compare(password, user.password);
    if (!match) {
      return res.status(401).json({ success: false, message: "Invalid credentials" });
    }

    // generateToken + res.cookie happen here (see Chapters 18-19)
    res.status(200).json({ success: true, data: { id: user._id, email: user.email } });
  } catch (err) {
    next(err);
  }
};
\end{lstlisting}

\begin{warnbox}[WARNING]
Return the same generic ``Invalid credentials'' message whether the email
does not exist or the password is wrong. Distinguishing the two tells an
attacker which emails are registered in your system.
\end{warnbox}

\section{Alternative: Hashing in a Mongoose Pre-Save Hook}

Instead of hashing manually in every controller that creates or updates a
user, you can centralize it in the User schema using a \texttt{pre("save")}
hook. This guarantees hashing happens no matter where a user document is
saved from.

\begin{lstlisting}[language=JavaScript]
// models/User.js
import mongoose from "mongoose";
import bcrypt from "bcrypt";

const userSchema = new mongoose.Schema({
  name: { type: String, required: true },
  email: { type: String, required: true, unique: true },
  password: { type: String, required: true },
  avatar: { type: String, default: "" },
  role: { type: String, enum: ["user", "manager", "admin"], default: "user" },
});

userSchema.pre("save", async function (next) {
  if (!this.isModified("password")) return next(); // skip if unchanged

  const salt = await bcrypt.genSalt(10);
  this.password = await bcrypt.hash(this.password, salt);
  next();
});

userSchema.methods.comparePassword = function (candidate) {
  return bcrypt.compare(candidate, this.password);
};

export default mongoose.model("User", userSchema);
\end{lstlisting}

\subsection{Why \texttt{isModified("password")} Matters}

Without this check, every time a user document is saved -- even to update an
unrelated field like \texttt{avatar} -- the hook would re-hash the
\emph{already-hashed} password. That produces a hash-of-a-hash, permanently
locking the user out because future logins compare against the wrong value.
\texttt{isModified("password")} ensures the hook only re-hashes when the raw
password field was actually changed on this save.

\begin{tipbox}[TIP]
With the pre-save hook approach, the controller simplifies to
\texttt{User.create(\{ name, email, password, role \})} -- no manual
\texttt{bcrypt.hash} call needed. Login still uses
\texttt{user.comparePassword(password)} (or \texttt{bcrypt.compare}
directly) exactly as shown above.
\end{tipbox}

\whatwhyhow
{bcrypt hashes passwords with a per-password random salt before storage, and compares candidate passwords against the stored hash on login.}
{Plain-text or reversibly-encrypted passwords turn any database leak into a full account compromise; one-way hashing with salt makes stored values useless to an attacker even if stolen.}
{Hash with \texttt{bcrypt.hash(password, 10)} on registration (or via a \texttt{pre("save")} hook guarded by \texttt{isModified("password")}), and verify with \texttt{bcrypt.compare(password, user.password)} on login.}
